\section*{Bab \ref{ch:b2:alpine-orogen} --- Anatomi Sabuk Pegunungan: Alpen}

\begin{solution}{exo:b2:alpine-orogen:1}
(a) $e = 1 - 27.0/36.0 = 0.25$, yaitu \qty{25}{\percent}; $S = 36.0 - 27.0 =
\qty{9.0}{km}$. (b) Lapisan itu harus mempertahankan panjangnya (tanpa peregangan atau penipisan di sepanjangnya),
penampang harus sejajar arah transpor, dan garis-garis pin harus menandai tempat
lapisan-lapisan tidak saling tergelincir.
\end{solution}

\begin{solution}{exo:b2:alpine-orogen:2}
(a) \omterm{def:b2:alpine-orogen:klippe}{Klippe}: gneis itu adalah sisa terpencil suatu \omterm{def:b2:alpine-orogen:nappes}{nappe} yang bertumpu di atas \omterm{def:b2:alpine-orogen:nappes}{autoktonnya};
erosi telah menyingkirkan sisa lembaran di sekelilingnya. (b) \omterm{def:b2:alpine-orogen:klippe}{Jendela tektonik}: erosi telah
memotong \omterm{def:b2:alpine-orogen:nappes}{nappe} gneis dan menyingkapkan batuan yang lebih muda di bawah sesar dorong basalnya, yang
karena itu miring menjauhi jendela di setiap sisi.
\end{solution}

\begin{solution}{exo:b2:alpine-orogen:3}
(a) Kompleks ultramafik (mantel atas); kompleks gabroik (kerak samudra bawah);
kompleks \omterm{def:b2:magma-emplacement:bodies}{retas} berlembar (kerak tengah); kompleks vulkanik, berstruktur bantal (kerak atas).
(b) Kompleks \omterm{def:b2:magma-emplacement:bodies}{retas} berlembar. (c) Sebagian besar kerak samudra-samudra luas habis oleh subduksi;
\omterm{def:b2:alpine-orogen:ophiolite}{ofiolit} yang bertahan di daratan sebagian besar terbentuk di cekungan samudra kecil atau di antara busur
kepulauan, sedangkan \omterm{def:b2:alpine-orogen:ophiolite}{ofiolit} Tethys terbentuk di tempat sebuah benua menipis sampai sebuah samudra kecil membuka.
\end{solution}

\begin{solution}{exo:b2:alpine-orogen:4}
(a) \omterm{def:b2:alpine-orogen:nappes}{Nappe} Monte Rosa dan Grand St Bernard: sisi Eropa (Briançonnais dan
tepi pasif Eropa); Zermatt--Saas dan zona Combin: samudra Piemont--Liguria; sistem Dent
Blanche: tepi Apulia. (b) Apulia, yang berakhir paling atas. (c) Lempeng atas
tidak tersubduksi, dan irisan-irisan yang dikerok dari lempeng bawah tertumpuk di bawahnya menurut urutan
kedatangan lempeng bawah: samudra jauh lebih dulu, tepi dekat terakhir dan paling bawah. Pelipatan tumpukan yang terjadi kemudian,
atau sesar dorong yang keluar dari urutan, dapat mengacaukan urutan itu.
\end{solution}

\begin{solution}{exo:b2:alpine-orogen:5}
(a) $2\pi G\,\Delta\rho = 2\pi\times6.674\times10^{-11}\times500 =
\qty{2.097e-7}{s^{-2}}$; $h = 1.70\times10^{-3}/2.097\times10^{-7} = \qty{8.1}{km}$.
(b) $2.097\times10^{-7}\times25\,000 = \qty{5.24e-3}{m.s^{-2}}$, sekitar
\qty{-524}{mGal}. (c) Akar itu bukan lempeng tak hingga: akar itu adalah tubuh sedalam beberapa puluh kilometer
dan lebarnya kira-kira hanya selebar sabuk, sehingga tarikannya pada suatu titik di permukaan lebih kecil daripada
lempeng setebal itu yang membentang tak hingga; tubuh rapat seperti zona Ivrea
menambahkan anomali positif, dan kontras massa jenis di kedalaman mungkin lebih kecil daripada yang diandaikan.
\end{solution}

\begin{solution}{exo:b2:alpine-orogen:6}
(a) $164 - 44 = \qty{120}{Ma}$. (b) $z = P/(\bar\rho g)$:
$2.8\times10^9/(3000\times9.8) = \qty{95.2}{km}$; $2.8\times10^9/(3300\times9.8) =
\qty{86.6}{km}$. (c) ``Lebih dari \qty{90}{km}'' terletak di dalam rentang ini: kedalamannya bergantung pada
massa jenis rata-rata yang dipilih untuk batuan di atasnya sama besarnya dengan pada tekanannya.
\end{solution}

\begin{solution}{exo:b2:alpine-orogen:7}
(a) Sayap sepanjang $L$ yang miring dengan sudut $\theta$ membentang $L\cos\theta$ secara horizontal; semua sayap
memiliki kemiringan yang sama, sehingga $l = l_0\cos\theta$ dan $e = 1 - l/l_0 = 1 - \cos\theta$.
(b) $\cos\theta = 0.70$: $\theta = \qty{45.6}{\degree}$. (c) $1 - \cos\qty{40}{\degree} =
\qty{23.4}{\percent}$, seperti pada gambar lipatan chevron; $1 - \cos\qty{60}{\degree}
= \qty{50}{\percent}$.
\end{solution}

\begin{solution}{exo:b2:alpine-orogen:8}
(a) $l_0 = 118/2.0 = \qty{59.0}{km}$; $S = 59.0 - 42 = \qty{17.0}{km}$; $e = 17.0/59.0 =
\qty{28.8}{\percent}$. (b) $A_e = 118 - 42\times2.0 = \qty{34.0}{km^2}$; $d = 34.0/17.0 =
\qty{2.0}{km}$. (c) Ya: \omterm{def:b2:alpine-orogen:decollement}{décollement} terletak di dasar penutup setebal \qty{2.0}{km}, di bawah
puncaknya; penutup itu terlepas pada dasarnya.
\end{solution}

\begin{solution}{exo:b2:alpine-orogen:9}
(a) Akar Airy adalah $4.5h$: pada \qty{46.5}{\degree} LU, $4.5\times2.74 = \qty{12.3}{km}$, yaitu
Moho pada \qty{42.3}{km}; pada \qty{45.55}{\degree} LU, $4.5\times0.65 = \qty{2.9}{km}$, yaitu Moho
pada \qty{32.9}{km}. (b) Di bawah daerah tinggi, Moho \qty{6.0}{km} lebih dangkal daripada ramalan
Airy; di bawah tepi Dataran Po, Moho \qty{28.3}{km} lebih dalam. (c) Akar itu tidak
terkompensasi secara lokal: akar itu bergeser ke selatan dan lebih besar daripada yang dituntut relief.
Lempeng melengkung sebagai satu kesatuan di bawah beban-beban di atasnya (penopangan regional yang dibahas dalam salah satu bab volume
Tahun 3), dan salah satu beban itu, tubuh Ivrea, sama sekali bukan topografi.
\end{solution}

\begin{solution}{exo:b2:alpine-orogen:10}
(a) $5.1/11.3 = 0.45$: kurang dari separuh. (b) $5.1 + 6.2 = \qty{11.3}{km}$. Puncak
Jura terletak di atas bidang lepasan atas; \qty{6.2}{km} yang terbawa di sepanjang bidang lepasan itu
lewat di bawah lipatan-lipatan penampang dan terpakai di struktur-struktur lebih jauh ke utara, sehingga
lapisan atas tidak pernah memendek sebesar itu di dalam penampang. (c) Satu lapisan penanda hanya merekam
pemendekan lapisan-lapisan tempat ia berada; setiap bidang lepasan di atas atau di bawahnya dapat
membawa pergeseran masuk ke atau keluar dari penampang. Seimbangkan beberapa lapisan, dan periksa dengan luas.
\end{solution}

\begin{solution}{exo:b2:alpine-orogen:11}
(a) Untuk pergeseran \qty{100}{km}, $\sqrt{200^2 + 100^2} = \qty{223.6}{km}$; untuk
\qty{50}{km}, $\sqrt{200^2 + 50^2} = \qty{206.2}{km}$; untuk \qty{150}{km},
$\sqrt{200^2 + 150^2} = \qty{250.0}{km}$. (b) $\arctan(100/200) = \qty{26.6}{\degree}$ dari
arah utara--selatan, garis konvergensi yang dekat dengan U\qty{150}{\degree}T yang dikutip Rosenberg dan
kawan-kawan. (c) Pergeseran itu tidak pasti sebesar $\pm\qty{50}{km}$, yang menggeser konvergensi
total sekitar \qty{20}{km} ke kedua arah; angka ketiga hasilnya jauh di bawah
ketidakpastian itu.
\end{solution}

\begin{solution}{exo:b2:alpine-orogen:12}
(a) $210/12.9 = \qty{16.3}{Ma}$ sampai $210/8.5 = \qty{24.7}{Ma}$ (\unit{km} dibagi
\unit{mm.yr^{-1}} menghasilkan juta tahun). (b) Lebih dari $4\times60 = \qty{240}{km}$; sekurang-kurangnya
$240 - 60 = \qty{180}{km}$ cekungan itu kini berada di bawah sesar dorong Alpen atau telah
tererosi. (c) Cekungan itu terbentuk di depan baji sesar dorong dan bermigrasi bersamanya: bagian
dalamnya selalu tertindih oleh sesar dorong yang membebaninya, dan isiannya yang tertua terperangkap
dalam sesar dorong, seperti \omterm{def:b2:alpine-orogen:foreland}{Molasse} Subalpin.
\end{solution}

\begin{solution}{pb:b2:alpine-orogen:1}
\noindent\textbf{Bagian I.}

\textbf{1.} $3.0/\cos\qty{35}{\degree} = \qty{3.66}{km}$; $2.5/\cos\qty{50}{\degree} =
\qty{3.89}{km}$; $4.0/\cos\qty{40}{\degree} = \qty{5.22}{km}$.

\textbf{2.} $8.00 + 2\times3.66 + 5.00 + 2\times3.89 + 6.0/\cos\qty{10}{\degree} +
2\times5.22 + 10.0/\cos\qty{5}{\degree} = 8.00 + 7.32 + 5.00 + 7.78 + 6.09 + 10.44 + 10.04 =
\qty{54.67}{km}$; dengan bentangan yang terulang, $l_0 = \qty{60.7}{km}$.

\textbf{3.} Lapisan di bawah lembaran sesar dorong adalah lapisan yang sama, yang terbawa ke bawah lembaran itu oleh
sesar dorong: sebelum pensesaran dorong, lapisan itu segaris dengan sisanya. Mengabaikannya akan memulihkan
lapisan itu terlalu pendek.

\textbf{4.} $e = 1 - 48.0/60.7 = 0.209$, yaitu \qty{20.9}{\percent}; $S = 60.7 - 48.0 =
\qty{12.7}{km}$.

\textbf{5.} $60.7 - 60.2 = \qty{0.5}{km}$, $0.5/60.7 = \qty{0.8}{\percent}$: dalam batas
ketelitian pengukuran; penampang itu dapat diterima.

\textbf{6.} $60.7 - 54.0 = \qty{6.7}{km}$, \qty{11}{\percent}: dua lapisan dari satu lembaran
tidak mungkin pulih ke panjang yang sejauh itu bedanya. Penampang itu tidak dapat diterima sebagaimana digambar: ada
struktur yang hilang atau salah gambar, garis pin di tempat yang salah, gerakan keluar dari bidang
penampang, atau bidang lepasan di antara kedua lapisan itu.

\noindent\textbf{Bagian II.}

\textbf{7.} $146/2.4 = \qty{60.8}{km}$.

\textbf{8.} Hasil itu berbeda \qty{0.1}{km} dari \qty{60.7}{km}: panjang garis dan luas sesuai.

\textbf{9.} $A_e = 146 - 48.0\times2.4 = 146 - 115.2 = \qty{30.8}{km^2}$.

\textbf{10.} $d = 30.8/12.7 = \qty{2.43}{km}$, sama dengan \qty{2.4}{km} tebal penutup dalam
batas ketelitian data: \omterm{def:b2:alpine-orogen:decollement}{décollement} adalah dasar penutup itu.

\textbf{11.} $160/2.4 = \qty{66.7}{km}$, \qty{6.0}{km} lebih panjang daripada panjang garis. Penebalan
evaporit di dasar penutup, yang mengalir ke dalam inti antiklin, menambah luas
tanpa menambah panjang; neraca luas dengan tebal tetap lalu gagal.

\textbf{12.} Penutup itu bergerak di atas \omterm{def:b2:alpine-orogen:decollement}{décollement} di dasarnya dan batuan dasar tidak
terlibat. Di bawah sabuk itu, puncak batuan dasar seharusnya datar atau miring landai dan
merata, tidak terputus oleh sesar dorong di atasnya.

\noindent\textbf{Bagian III.}

\textbf{13.} $22 + 18 + 25 = \qty{65}{km}$; $S = 65 - 15 = \qty{50}{km}$.

\textbf{14.} $e = 50/65 = \qty{76.9}{\percent}$, dibandingkan dengan \qty{20.9}{\percent} pada
sabuk daratan depan: \omterm{def:b2:alpine-orogen:nappes}{nappe} jauh lebih memendek.

\textbf{15.} $30 - 22 = \qty{8}{km}$ ($e = \qty{26.7}{\percent}$). Batuan dasar itu sendiri
memendek: \omterm{def:b2:alpine-orogen:skin}{tektonik kulit tebal}.

\textbf{16.} Sabuk daratan depan $12.7/(11 - 4) = \qty{1.81}{km/Ma}$, yaitu
\qty{1.81}{mm} per tahun; \omterm{def:b2:alpine-orogen:nappes}{nappe} $50/(30 - 20) = \qty{5.0}{mm}$ per tahun.

\textbf{17.} \omterm{def:b2:alpine-orogen:nappes}{Nappe} (30 sampai \qty{20}{Ma}) terdeformasi sebelum sabuk daratan depan (11 sampai
\qty{4}{Ma}): deformasi bermigrasi ke luar, ke arah daratan depan. Inilah urutan di
Alpen: sesar dorong Glarus pada Oligosen dan Miosen, masif eksternal pada Miosen,
Jura pada Miosen akhir dan Pliosen.

\textbf{18.} $12.7 + 50 + 8 = \qty{70.7}{km}$.

\noindent\textbf{Bagian IV.}

\textbf{19.} $\qty{0.5}{cm.yr^{-1}}\times\qty{40}{Ma} = 0.005\times4\times10^7 =
\qty{2.0e5}{m} = \qty{200}{km}$.

\textbf{20.} $200 + 18 = \qty{218}{km}$.

\textbf{21.} Hanya yang terawetkan atau tersimpulkan yang dihitung: sebagian samudra mungkin telah
tersubduksi sebelum \qty{85}{Ma} atau lebih cepat daripada laju rata-rata, irisan \omterm{def:b2:alpine-orogen:ophiolite}{ofiolit} mungkin telah
tererosi atau terbawa keluar dari penampang, dan konvergensi yang miring terhadap penampang membawa
masuk dasar samudra yang tidak terlihat oleh penampang.

\textbf{22.} Zona internal $60 - 20 = \qty{40}{km}$ ($e = \qty{66.7}{\percent}$); baji belakang
$70 - 48 = \qty{22}{km}$ ($e = \qty{31.4}{\percent}$).

\textbf{23.} $12.7 + 50 + 8 + 40 + 22 = \qty{132.7}{km}$; berada dalam rentang taksiran
untuk Alpen Tengah, sekitar 125 sampai \qty{160}{km} pemendekan tumbukan (sebelum ditambah
gelinciran mendatar).

\textbf{24.} Lebar pulihan $153 + 132.7 = \qty{285.7}{km}$; kerak pulihan
$285.7\times30 = \qty{8571}{km^2}$; kerak masa kini $153\times45 = \qty{6885}{km^2}$. Yang hilang:
$8571 - 6885 = \qty{1686}{km^2}$, \qty{19.7}{\percent} kerak pulihan; kerak itu telah
tererosi (dan tercurah ke \omterm{def:b2:alpine-orogen:foreland}{cekungan depan-pegunungan}) atau tersubduksi bersama kerak bawah ke dalam
mantel. Tanpa kehilangan itu, kerak akan setebal rata-rata $8571/153 = \qty{56}{km}$.
Pemendekan tumbukan total di sepanjang penampang adalah
\textbf{\qty[text-series-to-math=true]{133}{km}}.
\end{solution}
